
L'idéal engendré par @@M0@@ est un @@M1@@-idéal. Pour le voir, on remarque d'abord que les relations @@M2@@ pour @@M3@@ entraînent que les @@M4@@ pour @@M5@@ sont multiples de @@M6@@. Si on se place dans @@M7@@, l'image de @@M8@@ est un polynôme de degré @@M9@@, et un argument universel montre alors que ses @@M10@@ pour @@M11@@ sont égaux à @@M12@@ (cf. V 4.8). Comme dans @@M13@@ on a @@M14@@, on en déduit que les @@M15@@ pour @@M16@@ sont multiples de @@M17@@. D'après les relations V 2.4.2, cela entraîne que l'idéal engendré par @@M18@@ est un @@M19@@-idéal.

On pose alors @@M20@@. L'homomorphisme naturel @@M21@@ sera noté @@M22@@; c'est un @@M23@@-homomorphisme. D'autre part, @@M24@@ est un @@M25@@-module libre ayant pour base les éléments @@M26@@. Il résulte de la relation @@M27@@ et de l'hypothèse d'inversibilité faite sur @@M28@@ que @@M29@@ est inversible. Par suite @@M30@@ admet pour base sur @@M31@@ les éléments @@M32@@ ; on définit alors une application @@M33@@ en associant à tout élément @@M34@@ de @@M35@@ l'élément @@M36@@ de @@M37@@. Il est clair que l'application ainsi définie est un homomorphisme additif, @@M38@@-linéaire, et tel que @@M39@@.

On munit ensuite @@M40@@ d'une structure de @@M41@@-anneau, en posant :
@@M42@@
Dans ces formules, @@M43@@ est considéré comme @@M44@@-algèbre par l'homomorphisme @@M45@@.
On pose ensuite, pour @@M46@@ et @@M47@@ :
@@M48@@
